\chapter{Proposal Feedback and Iteration Plan}
% ============================================================

\section{Proposal Feedback Summary}
\label{sec:feedback-summary}

The proposal defence panel returned six comments. Two concerns were each raised
twice, which the team treated as indicating priority.

\begin{enumerate}
  \item The proof of concept was assessed as crude, and the panel separately
        reiterated that the team needs to work on the PoC. This was the most emphasised
        item.
  \item The evaluation matrix was assessed as requiring greater precision, and the
        panel separately noted that more effort should be made in the context of the
        evaluation metric.
  \item The team was unable to explain the distinction between static and dynamic
        environments when questioned.
  \item More in-depth background study was required.
  \item The panel suggested better simulators, and add-on applications where the
        outcomes could be utilised.
  \item The application perspective needed to be made stronger.
\end{enumerate}

The panel's overall assessment recorded the project as a good niche project with a
good direction, subject to the above.

\section{Proposal Action Register}
\label{sec:action-register}

\begin{longtable}{p{0.04\textwidth} p{0.22\textwidth} p{0.34\textwidth} p{0.10\textwidth} p{0.20\textwidth}}
\caption{Proposal action register}\\
\toprule
\textbf{No.} & \textbf{Feedback / required change} & \textbf{Action taken} & \textbf{Status} & \textbf{Evidence} \\
\midrule
\endfirsthead
\toprule
\textbf{No.} & \textbf{Feedback / required change} & \textbf{Action taken} & \textbf{Status} & \textbf{Evidence} \\
\midrule
\endhead
1 & PoC is crude; students need to work on the PoC &
Replaced the proposal-stage illustrative mock-up with an executable simulation:
four-zone TSN topology with gPTP, three instrumented attack configurations, 15 seeded
runs, and a feature extraction pipeline producing 4{,}740 labelled windows. The
operator-facing surface is now explicitly labelled as a presentation mock-up and is no
longer presented as the proof of concept. &
Addressed &
Ch.~\ref{ch:iteration1}; Fig.~\ref{fig:topology}; Sec.~\ref{sec:demo-evidence}; Sec.~\ref{sec:test-results} \\
\addlinespace
2 & Evaluation matrix requires greater precision &
Replaced the three-row metric summary with a full evaluation specification: per-class
precision/recall/F1 and macro-F1, detection latency in sync intervals, false alarms per
vehicle-hour, worst-client macro-F1, centralized-to-FedAvg gap, convergence rounds,
update volume, and zone-attribution top-1 accuracy reported separately from detection.
Accuracy removed as a headline measure, with class-imbalance justification. &
Addressed &
Sec.~\ref{sec:success-criteria}; Sec.~\ref{sec:testing-strategy}; Sec.~\ref{sec:validation-plan} \\
\addlinespace
3 & Unable to explain static vs.\ dynamic environment &
Both senses are now stated explicitly as load-bearing assumptions: the network
environment is statically configured (A1), and training is static/one-shot rather than
continuous (A2). The dependence of the detection approach on a statically known
expected value is stated, and dynamic reconfiguration is recorded as an explicit
limitation. &
Addressed &
Sec.~\ref{sec:assumptions} (A1, A2); Sec.~\ref{sec:arch-decisions}; Sec.~\ref{sec:limitations} \\
\addlinespace
4 & More in-depth background study required &
Literature review extended beyond the five sources presented at defence, recorded
against methodology, results, limitations and resulting design decision; a
research-to-decision table now links each finding to a concrete design choice; a
comparison against the departmental FYP repository was added. &
Addressed &
Sec.~\ref{sec:research-method}--\ref{sec:limitations-gaps}; Sec.~\ref{sec:research-findings}; Table~\ref{tab:dims} \\
\addlinespace
5 & Better simulators; add-on applications where outcomes can be utilised &
Simulation framework and version justified against alternatives in an architecture
decision record. Consumption of detector output defined concretely: the vSOC incident
queue, and the in-vehicle integrity decision consumed by fault-tolerant timing
failover. &
Partially addressed &
Sec.~\ref{sec:arch-decisions}; Sec.~\ref{sec:stakeholders}; Sec.~\ref{sec:deployment-plan} \\
\addlinespace
6 & Application perspective needs to be stronger &
Added deployment positioning (zone controller and central vehicle compute), the
consuming stakeholders and their required evidence, and the deployment plan
identifying what the detector's output drives. &
Partially addressed &
Sec.~\ref{sec:problem-context}; Sec.~\ref{sec:stakeholders}; Sec.~\ref{sec:deployment-plan} \\
\bottomrule
\end{longtable}

% [[VERIFY]] Items 5 and 6 are marked "Partially addressed" deliberately. Claiming
% full closure on an application-perspective comment with no deployed consumer is
% not defensible. If you can show the alert surface driven by real model output
% before the evaluation, upgrade them.

\section{Iteration Strategy}
\label{sec:iteration-strategy}

Development follows four time-boxed iterations, each ending in a demonstrable
artefact rather than a document. The ordering is deliberately
\emph{environment $\rightarrow$ data $\rightarrow$ model $\rightarrow$ federation},
on the principle that no modelling decision can be validated before the data
generation path is trustworthy. An explicit decision gate separates Iteration~02 from
Iteration~03: the detection claim is tested centrally, on a single machine, before any
federated infrastructure is built, so that a null feature result is discovered before
it is built upon.

\section{Iteration 1 Plan}
\label{sec:iter1}

\textbf{Objective.} Establish a reproducible simulation environment in which TSN
timing attacks produce measurable, repeatable effects.

\begin{itemize}
  \item \textbf{Selected features:} zonal topology construction; gPTP synchronisation;
        attack injection harness; seeded reproducible execution.
  \item \textbf{Deliverables:} \texttt{CarazoneZonalNetwork.ned},
        \texttt{omnetpp.ini} with four named configurations, executable simulation.
  \item \textbf{Evidence:} topology rendering, benign vs.\ attack offset traces
        across multiple seeds.
  \item \textbf{Exit criterion:} an injected attack produces a timing deviation
        distinguishable from the benign baseline, reproducibly.
  \item \textbf{Status: complete.}
\end{itemize}

\section{Iteration 2 Plan}
\label{sec:iter2}

\textbf{Objective.} Convert raw simulation output into a labelled, windowed,
protocol-native feature dataset.

\begin{itemize}
  \item \textbf{Selected features:} passive signal capture; windowing; feature
        computation; labelling with zone and class attribution; dataset export.
  \item \textbf{Deliverables:} \texttt{extract\_features.py}, \texttt{features.csv}.
  \item \textbf{Evidence:} 225 vectors across 15 runs; 4{,}740 windows; class and
        per-zone distribution.
  \item \textbf{Exit criterion:} a labelled dataset regenerable from configuration and
        seed.
  \item \textbf{Status: complete.} Labelling is correctly attributed to the victim
        zone only; the target zone is rotated across runs so that every zone appears
        as both victim and bystander (Section~\ref{sec:test-results}).
\end{itemize}

\section{Work Breakdown and Responsibilities}
\label{sec:workbreakdown}

\begin{longtable}{p{0.2\textwidth} p{0.37\textwidth} p{0.37\textwidth}}
\caption{Work breakdown by member}\label{tab:workbreakdown}\\
\toprule
\textbf{Area} & \textbf{Adeena Zaib} & \textbf{Khaula Bilqees} \\
\midrule
\endfirsthead
\toprule
\textbf{Area} & \textbf{Adeena Zaib} & \textbf{Khaula Bilqees} \\
\midrule
\endhead
Setup &
OMNeT++/INET zonal topology; model architecture &
TSN-native signal instrumentation; federated learning environment setup \\
\addlinespace
Data generation &
Attack traffic simulation, all scenarios, per zone &
Feature specification review; labelling scheme \\
\addlinespace
Federated learning &
Per-zone dataset packaging for clients &
Federated client/server implementation \\
\addlinespace
Training and evaluation &
Attack-specific ground-truth verification &
Model training; hyperparameter tuning; metrics collection \\
\addlinespace
Shared &
Literature review; architecture diagrams; report writing &
Literature review; architecture diagrams; report writing \\
\bottomrule
\end{longtable}


Each member is responsible for being able to explain, modify and defend the complete
architecture, not only their primary component. No member's ownership is limited to
documentation, presentation or manual testing.

\section{Progress Against Plan}
\label{sec:progress}

\begin{longtable}{p{0.16\textwidth} p{0.38\textwidth} p{0.14\textwidth} p{0.24\textwidth}}
\caption{Planned milestones versus actual progress}\\
\toprule
\textbf{Milestone} & \textbf{Planned outcome} & \textbf{Status} & \textbf{Note} \\
\midrule
\endfirsthead
\toprule
\textbf{Milestone} & \textbf{Planned outcome} & \textbf{Status} & \textbf{Note} \\
\midrule
\endhead
Iteration 01 & Simulation environment; attacks instrumented & Complete & Four named
configurations implemented \\
Iteration 02 & Traffic generation; feature pipeline; labelled per-zone datasets &
Complete & Cross-zone residual feature outstanding \\
Iteration 03 & Federated learning pipeline; initial training runs & Not started &
Preceded by residual feature and zone heterogeneity (Section~\ref{sec:delays}) \\
Iteration 04 & One-shot federated training; evaluation metrics collected & Not started & n/a \\
\bottomrule
\end{longtable}

\section{Delays, Deviations and Recovery Plan}
\label{sec:delays}

\textbf{D1. Detection model not yet implemented.} The feature dataset is currently the
end of the implemented pipeline. \emph{Cause:} Iterations 01--02 absorbed more effort
than planned, primarily in instrumenting protocol-internal signals that are not
exposed by default. \emph{Recovery:} a centralized detector is scheduled as the first
Iteration~03 task, deliberately ahead of any federated infrastructure, so that the
feature claim (SC-2) is tested before further work depends on it.

\textbf{D2. Zone symmetry.} All four zones of the Iteration~1 topology carry identical
traffic and produce an identical number of feature windows. \emph{Consequence:}
federated, centralized and local-only training would converge to equivalent models,
making the federated comparison (SC-4) null by construction. \emph{Recovery:} the
five-zone reference topology (Section~\ref{sec:ref-topology}) builds asymmetry into the structure itself,
with between three and seven endpoint ECUs per zone and distinct traffic classes by
zone function, giving quantity, feature and label skew by design. Differentiated
per-zone traffic profiles and the measurement of the resulting skew remain Iteration~03
tasks, before any federated training run.

\textbf{D3. Cross-zone residual feature added after mid-iteration review.} The
feature set initially computed statistics within each zone in isolation, with no
comparison between a zone and its siblings. \emph{Consequence:} the implemented
features were timing statistics that a generic detector could also compute, which did
not support the protocol-native claim SC-2 is designed to test. \emph{Resolution:} the
extraction loop was restructured so that the observation window is the outer iteration,
making every zone's controller offset available for the same window, and two residual
features were added. The ablation in Section~\ref{sec:acceptance-validation} confirms these carry the detection
signal. \emph{Closed.}

\textbf{D4. Deviations from the approved proposal.} Three deviations are recorded:
the number of federated clients changed from three to four and then to five, in each
case following the topology
design; team composition is two members rather than the three assumed by the
proposal's ownership model, with responsibilities reallocated accordingly
(Section~\ref{sec:workbreakdown}); and the attack surface was broadened to include FRER (IEEE 802.1CB)
stream-redundancy manipulation alongside gPTP synchronisation attacks, which the
approved proposal did not name. FRER remains committed scope and is in progress; it
is not yet represented in the Iteration~1 dataset because it requires a redundant
topology, and that extension is scheduled in Section~\ref{sec:remaining-iterations}.

% ============================================================